\section{\sestina{} architecture}
\label{sec:agent}

\sestina{} is a configuration of the \vsix{} policy family rather than a new architecture. This
section describes the agent conceptually: Figure~\ref{fig:pipeline} shows the decision
pipeline, and Table~\ref{tab:components} sets its components against \vsix{}'s. The literal
configuration string, its option map, the parameter vector layout and the freeze artifact's hashes
are in Appendix~\ref{app:parameters}.

\begin{figure}[!ht]\centering
\begin{tikzpicture}[
  font=\small,
  box/.style={draw=fishgray,rounded corners=2pt,align=center,inner sep=5pt,minimum height=9mm},
  % the three pipeline stages carry the same amount of content, so they are
  % held to one height rather than left to differ by the 1pt their last lines do
  stage/.style={box,minimum height=15mm,text width=32mm},
  pub/.style={box,fill=fishlight,draw=fishgreen},
  arr/.style={-{Stealth[length=2mm]},fishgray,thick}]
\node[pub] (hist) at (0,1.15) {public event stream\\{\scriptsize asks, answers, transfers,}\\[-1pt]{\scriptsize declarations, card counts}};
\node[box] (hand) at (0,-1.05) {own hand\\{\scriptsize nine cards}};
\node[stage] (bel)  at (4.3,0.05) {approximate deal belief\\{\scriptsize Sinkhorn fit,}\\[-1pt]{\scriptsize behavioural weighting}};
\node[stage] (pol)  at (8.2,0.05) {linear ask and\\declaration score\\{\scriptsize + public-history tie-break}};
\node[stage] (srch) at (12.1,0.05) {guarded search\\{\scriptsize endgame only,}\\[-1pt]{\scriptsize paired LCB deviation}};
\node[align=center] (act) at (12.1,-1.75) {\textbf{action}};
\draw[arr] (hist) -- (bel); \draw[arr] (hand) -- (bel);
\draw[arr] (bel) -- (pol); \draw[arr] (pol) -- (srch); \draw[arr] (srch) -- (act);
\draw[arr,dashed] (pol.south) .. controls +(0,-1.2) and +(-2,-1.2) .. (act.west);
\node[font=\scriptsize,fishgray,align=center] at (10.15,-1.25) {blueprint action\\if search declines};
% The enclosure is fitted to the two nodes it contains rather than drawn at
% hand-measured coordinates: at (-1.75,1.75) its sides ran along the borders of
% the public-event-stream box, which is exactly as wide as the rectangle was.
\node[draw=fishgreen,dashed,rounded corners=3pt,inner sep=6pt,fit=(hist)(hand)] (iset) {};
\node[font=\scriptsize,fishgreen,below=2pt of iset] {information set};
\end{tikzpicture}
\caption{The information structure of a Fish seat and the \sestina{} decision pipeline. A seat
observes the public event stream and its own hand, and nothing else; there is no path from a policy to
the deal or to another seat's hand (\S\ref{sec:inference-engine}). Because every card movement is
public, that information set determines an exact constraint-conditioned distribution over the
initial deal (\S\ref{sec:inference}). The deployed path does not construct it: it runs approximate
constraint-based inference with behavioural weighting, a Sinkhorn fit that starts from the
ask-history weighting and imposes the same constraints iteratively (\S\ref{sec:inference-prior}).
The first two stages together, the belief and the linear ask and declaration score with its
tie-break, form the \emph{blueprint}, the non-searching policy. The search runs only in the endgame
and departs from the blueprint's choice only when the guarded selection rule of
\S\ref{sec:agent-search} favours the departure.}
\label{fig:pipeline}
\end{figure}

\begin{table}[!ht]\centering
\caption{Components of \sestina{} against \vsix{}. Harness identifiers appear in parentheses and
are used in the paper only where a command must be reproduced.}
\label{tab:components}
{\tabfont\input{tables_v07/components}}
\end{table}

\subsection{Action space and ask policy}
\label{sec:agent-policy}

At its turn a seat faces one decision with two branches: declare a half-suit, or ask. Declaration is
treated first in the code and is described in \S\ref{sec:agent-declare}; this subsection covers the
ask.

\paragraph{The legal action set.} With $\mathcal{O}_s$ the opponents and $\mathcal{H}_t$ the
half-suits not yet claimed, the rules of \S\ref{sec:game-rules} give
\begin{equation}
  \mathcal{A}_t \;=\; \bigl\{(c,q) \;:\; q \in \mathcal{O}_s,\; \text{hand}_t(q) \neq \emptyset,\;
  H(c) \in \mathcal{H}_t,\; c \notin \text{hand}_t(s),\;
  \text{hand}_t(s) \cap H(c) \neq \emptyset \bigr\}.
  \label{eq:legal}
\end{equation}
The last condition encodes the ownership condition that makes an ask informative: the asker must
already hold a card of the half-suit, so asking publishes partial ownership.

What the policy actually enumerates is a strict subset of \eqref{eq:legal}. A \emph{live-ask filter} first removes candidates the seat can already prove are misses, that is,
those whose target is excluded by a certificate, $q \notin \mathcal{M}(c)$. It is this filter, rather
than any weighting inside \eqref{eq:askscore}, that keeps a provably unsuccessful ask off the score. It
binds: it removes what would otherwise have been the argmax at roughly \vsevenGateBindRate\% of
\vsevenGateBindN{} ask decisions captured from \sestina{}'s own seats. If no candidate survives the filter the full legal set is
restored, so the guarantee is not absolute.

\paragraph{The score.} Each candidate is scored by a fitted linear function of a
$20$-dimensional feature vector $\phi(a)$, plus a learned value term and, in \sestina{}, three
additional ownership terms:
\begin{equation}
  u(a) \;=\; \lambda\, w^{\!\top}\phi(a) \;+\; \beta\, V_{\mathrm{ask}}(a)
  \;+\; \lambda\bigl(w_{\mathrm{void}}\, g_1(a) + w_{\mathrm{team}}\, g_2(a)
  + w_{\mathrm{last}}\, g_3(a)\bigr),
  \label{eq:askscore}
\end{equation}
with $\lambda = \vsevenLinearWeight{}$ and $\beta$ scale parameters of the fit. \sestina{} adds a
fourth block of \vsevenResponderTotal{} extended coordinates, of which
\vsevenResponderNonzero{} is non-zero; that one is \eqref{eq:contest} and
\S\ref{sec:agent-contest} gives it. Writing $p_a = \hat\mu(c,q)$ for the modelled
hit probability, the features $\phi$ are functions of the belief and the public state. They include $p_a$ and $p_a^2$; an indicator that the card is publicly
certain; the acting seat's holding in $H(c)$ and the posterior mass its team already carries; the
best continuation probability elsewhere in the half-suit; the
reply threat and exposure a miss would create, both weighted by $1-p_a$; the target's hand size and
an indicator that this ask would empty it; the binary entropy $\mathrm{H}_2(p_a)$; and a score-margin
term. Appendix~\ref{app:parameters} lists all twenty with their fitted weights.

The three extra terms are the ones \vsix{} added, and each expresses an effect the twenty could not:
\begin{align}
  g_1(a) &= p_a \prod_{c' \in H(c) \setminus \{c\}} \bigl(1 - \hat\mu(c',q)\bigr),
    \label{eq:gvoid}\\
  g_2(a) &= \Pr\bigl[\,\mathrm{own}(c) \in \mathcal{T}_s \cup \{s\}\,\bigr],
    \label{eq:gteam}\\
  g_3(a) &= p_a \cdot \ind\bigl[\,q \text{ is the last opponent holding cards, and holds one}\,\bigr].
    \label{eq:glast}
\end{align}
Equation~\eqref{eq:gvoid} is an independence-based proxy for the probability that the ask takes
the target's \emph{last} card of the half-suit: it multiplies the marginals as if the target's
holdings of the other cards were independent, which the count-constrained distribution does not
imply. Equation~\eqref{eq:gteam} is the team-ownership probability of the named card, the belief
mass the acting team already carries, and \eqref{eq:glast} is an indicator that the ask would
trigger the forced endgame. Equation~\eqref{eq:gvoid} expresses the loss of the target's access to
that half-suit: because a legal ask requires holding a card of the half-suit, removing a target's
last card of $H(c)$ permanently removes that target's right to ask in it.

\paragraph{Selection.} The policy takes the argmax, with the tie handling of
\S\ref{sec:agent-tiebreak} and, when enabled, the search of \S\ref{sec:agent-search}:
\begin{equation}
  a^\star \;=\; \operatorname*{arg\,max}_{a \in \mathcal{A}_t} u(a).
  \label{eq:argmax}
\end{equation}

\sestina{} inherits $w$, $\lambda$, $\beta$ and the three extra weights from \vsix{} unchanged.
No gradient or population fit was run in this cycle: the one coordinate that is not
inherited, \eqref{eq:contest}, had its value chosen as the argmax of a measured dose response, which
is a one-dimensional selection on win rate and is reported as such in
\S\ref{sec:results-selection}. The vector was found by \vsix{}'s repaired
optimiser under a minimax-regret objective over a four-member opponent panel. Of the
\vsevenCoordinates{} coordinates in the extended vector a number are known to be inert, and
Appendix~\ref{app:parameters} lists which and why.

\subsection{Public-history tie-breaking}
\label{sec:agent-tiebreak}

In the \vsix{} mirror, \vsevenTieContested\% of contested ask decisions end with two or more
candidates scoring bit-for-bit identically under \eqref{eq:askscore}
(\S\ref{sec:background-vsix}), and the winner was whichever the enumerator emitted first, under an
unstable sort. That is a specification problem rather than a strength problem. The choice is
arbitrary, and it is determined by an implementation detail: the result of an unstable sort is
reproducible given the same implementation and inputs, but the policy does not specify it, so a
teammate modelling the seat would have to model the sort rather than the rule. The defensible
benefit of the rule below is a specified, reproducible selection rule that reads only the public
record.

Order the candidates by score and let the tied set be
\begin{equation}
  \mathcal{E}_t \;=\; \bigl\{a \in \mathcal{A}_t \;:\; u(a) \ge u(a_{(1)}) - \varepsilon \bigr\},
  \label{eq:tieset}
\end{equation}
where $a_{(1)}$ is the top-ranked candidate and $\varepsilon$ is a fixed tolerance. \sestina{}
replaces enumeration order with a deterministic draw keyed by the public record:
\begin{equation}
  a^\star \;=\; \mathcal{E}_t\bigl[\,\iota(\rho,\,|\mathcal{E}_t|)\,\bigr],
  \qquad
  \rho \;=\; \mathrm{hash}\bigl(\mathcal{P}_t,\; s,\; t\bigr),
  \label{eq:tiebreak}
\end{equation}
with $\mathcal{P}_t$ a partial digest of the public event stream, namely a rolling hash over each
event's kind, actor, target, card and outcome, which omits the half-suit index and the allocation a
declaration names; whether those omitted public fields can matter to a tie case was not
established. Here $s$ is the seat index, $t$ the event count, and $\iota$ a uniform index map from
the hash into the tied set. The tolerance
$\varepsilon$ is $10^{-9}$ in absolute score units, a fixed engine constant rather than a
configuration option, so $\mathcal{E}_t$ is a tolerance-based tied set rather than necessarily a
bit-identical one. Every seat observes all three arguments of $\rho$, so every teammate can
recompute it; no seat's private hand enters it. The index is applied to the tied block in the order
the score sort leaves it, and the arrangement of equal keys under that sort is fixed by the
implementation rather than by the rule, so what \eqref{eq:tiebreak} specifies is the index into the
block; an identical build reproduces the arrangement, as the isolation audit of
\S\ref{sec:results-controls} confirms for every non-searching decision.

Two scope conditions belong with that. \eqref{eq:tiebreak} governs the non-search path only: when
the endgame search fires it supplies its own selection (\S\ref{sec:agent-search}) and index $0$ is
whatever the sort left first. And the search draws its determinizations from a private per-seat
stream seeded from the harness, so the policy is a deterministic function of
$(\text{own hand}, \text{public record}, \text{reset seed})$ rather than of the first two alone. That
third argument is exactly what the side-channel audit of \S\ref{sec:evaluation-side} replays, which
is why the audit is a test of seat isolation and not of determinism.

The engine also offers a variant that draws $\rho$ from a private per-seat stream. That variant is
\emph{not} used: it would break the reconstructibility that makes \eqref{eq:tiebreak} a coordination
device. \S\ref{sec:development-tie} reports what happened when a later experiment tried to
\emph{select} within $\mathcal{E}_t$ rather than randomise over it.

\subsection{Half-suit contestation}
\label{sec:agent-contest}

The component with the largest observed leave-one-out drop in the holdout ablations
(\S\ref{sec:results-attribution}) is a single additional coordinate, one non-zero entry out of
\vsevenResponderTotal{} in an extended block of the parameter vector. Let $M_H$ be the expected number of cards of half-suit $H$ held by the
opposing team, and $q_H$ the number whose holder this seat cannot place. The term is
\begin{equation}
  r(a) \;=\; (1 - p_a)\cdot\frac{M_{H(c)}}{6}\cdot\frac{q_{H(c)}}{6},
  \label{eq:contest}
\end{equation}
entering \eqref{eq:askscore} with a \vsevenContestSign{} coefficient of
\texttt{\vsevenContestWeight{}}, so a large value \emph{encourages} the ask. It selects asks that are
likely to miss, into half-suits the opponents dominate, whose ownership is still ambiguous.

\paragraph{Feature interpretation.} The feature was designed to price \emph{information denial}: a
failed ask publishes a negative certificate, and the intent was that the fitted coefficient would
come out negative, meaning refuse to hand that certificate over. It came out positive. The negative
branch was swept separately and is flat to negative throughout, reaching \vsevenContestNegDose{}~pp
at the mirror-image dose, so the measured dose response does not support the design intent.

What the positive branch does was measured on the target rather than inferred. Over
$\sim$\vsevenContestDecisions{} target ask decisions, the target's ask accuracy falls from
\vsevenContestTargetAskBefore{} to \vsevenContestTargetAskAfter{}, while its declaration accuracy
\emph{rises} and it becomes \emph{less} often urgent, which is inconsistent with a declaration-channel
mechanism, where this phase expected to find the effect. The reading these measurements are
consistent with is contestation: asking into half-suits the opponents are trying to complete
degrades the quality of their asks, and the acting team pays part of the same cost, its own ask
accuracy falling from \vsevenContestSelfAskBefore{} to \vsevenContestSelfAskAfter{}. Three sibling
coordinates written to express the same denial idea were measured at the same dose and are inert or
harmful, which is consistent with a benefit specific to this feature combination in the tested
ablations rather than to opponent mass in general; a positive coefficient and the associated changes
in opponent accuracy suggest a mechanism without establishing it as the unique explanation.

Two facts about its provenance matter for \S\ref{sec:results-attribution}. It was found by a
coordinate sweep during adversary generation rather than by architecture work, and it is hand-set:
the value \texttt{\vsevenContestWeight{}} is the argmax of a dose response measured at
\vsevenPhaseTwoCell{} games a dose, not the output of a gradient or population fit. And it is not
uniformly good. Against the then-deployed policy it scored \vsevenOneSwitchArm{}~pp, but its worst
cell over the \vsix{} frontier is \vsevenOneSwitchArmWorst{}~pp against
\texttt{\vsevenOneSwitchArmWorstOpp{}}, which is negative, so it is a strength gain with a worst case
rather than an exploit (\S\ref{sec:development-adversaries}).

\subsection{Declaration policy}
\label{sec:agent-declare}

A declaration names, for each of the six cards of a half-suit, which member of the declaring team
holds it. It is all-or-nothing: correct, and the team scores the half-suit; wrong, and the opponents
do. The decision therefore has two parts, choosing what to claim and choosing when.

\paragraph{Allocation selection.} For half-suit $h$ the reference objective is the most probable
allocation among the acting team under the constraint-conditioned distribution,
\begin{equation}
  \alpha^\star(h) \;=\; \operatorname*{arg\,max}_{\alpha:\,H \to \mathcal{T}_s \cup \{s\}}
  \; \pi_{\mathrm{alloc}}(h,\alpha),
  \label{eq:bestalloc}
\end{equation}
with $\pi_{\mathrm{alloc}}$ from \eqref{eq:alloc}. Since \eqref{eq:alloc} depends on
$\alpha^\star$ only through its count signature, the argmax in \eqref{eq:bestalloc} is generally
not unique; the reference implementation takes the lexicographically first feasible assignment
among the maximisers. Equation~\eqref{eq:bestalloc} defines the reference objective; the deployed
approximation is described here. Three procedures must be kept
distinct. \emph{Marginal-product selection}: \sestina{} enumerates the capacity-feasible
assignments of the unresolved cards of $h$ to its own team and takes the argmax of the product of
per-card approximate marginals, $\prod_{c} \hat\mu(c,\alpha(c))$; that product is the selection
score only. \emph{Sequential approximate joint scoring}: the selected allocation is then priced
separately, by sequential conditioning, in which the belief is refit after each card of the
allocation is fixed in turn, and the resulting chain-rule product is what the policy reports as
$\hat\pi^\star_{\mathrm{alloc}}$ and thresholds on. \emph{Exact constraint-joint selection}: the
exact object \eqref{eq:alloc} is available from the block dynamic program but sits on no decision
path. \S\ref{sec:development-alloc} reports the relevant measurements: the marginal-product
selection and the sequential approximate joint's argmax selected the same allocation in every one
of \vsevenLoneDeclsTotal{} recorded voluntary declarations, of which \vsevenLoneDeclsDistinct{} are
from distinct (bank, urgency-setting) capture cells and the remainder from re-runs of the same
seeds and from the two cells of the sequential-joint policy, an agreement that is empirical and
does not establish equivalence to the exact joint maximiser; and substituting the exact constraint-joint selection
changed a small fraction of decisions and performed worse on the declaration decision.

\paragraph{Declaration timing.} Write $\hat\pi_{\mathrm{alloc}}^\star$ for the sequentially
estimated probability of the selected allocation and $\hat\pi_{\mathrm{team}}$ for the deployed
estimate of the probability the team owns $h$ outright, which is the larger of
$\hat\pi_{\mathrm{alloc}}^\star$ and the product over the half-suit's unresolved cards of the team's
summed marginals, $\prod_{c}\sum_{p \in \mathcal{T}_s \cup \{s\}}\hat\mu(c,p)$, an independence
product rather than \eqref{eq:team}. The value function is written on the
assumption that waiting is free, on the ground that an opponent cannot claim a half-suit the acting
team provably holds. That ground binds strictly only in the locked case, and the rule is applied
down to a team-ownership probability well below one, so the assumption is an approximation rather
than an identity. Under it the decision is an approximate comparison between immediate declaration
value and the value of waiting: declaring scores a point now with probability
$\hat\pi_{\mathrm{alloc}}^\star$, and publishes the team's holdings either way. The agent evaluates
that with its learned state-value function $V$:
\begin{equation}
  \Delta \;=\; \underbrace{\hat\pi_{\mathrm{alloc}}^\star\, V(\text{state} \mid \text{correct})
  \;+\; (1-\hat\pi_{\mathrm{alloc}}^\star)\, V(\text{state} \mid \text{wrong})}_{\text{declare}}
  \;-\; \underbrace{V(\text{state})}_{\text{wait}},
  \label{eq:declev}
\end{equation}
and declares when $\Delta > \delta$ for a fitted margin $\delta$. Both conditioned states differ from the waiting state in having the \emph{unresolved} cards of $h$
removed from the unresolved set, typically fewer than six, and in the deltas of the four
half-suit-control features the value function reads, all of which are functions of the team's
expected share of $h$ before the claim. They differ from \emph{each other} only in the score margin,
by $+1$ and $-1$.

A gate precedes \eqref{eq:declev}: a half-suit is priced by the value rule only if its deployed
team-ownership estimate $\hat\pi_{\mathrm{team}}$ clears a fitted floor, which is lowered under
pressure stage 1 and removed under stage 2. Two overrides sit above \eqref{eq:declev}. Under the
forced endgame, when the opposing team is cardless and the acting team must resolve the remaining
half-suits among itself, the engine polls the team down the willingness ladder of
\S\ref{sec:game-dialect}, and at each level the agent is willing to declare a half-suit when
$\hat\pi_{\mathrm{alloc}}^\star \ge \tau$ for that level's threshold $\tau$; the final level is
compulsory. Under \emph{pressure}, the escalation described next, stage 1 lowers the thresholds,
declaring when $\hat\pi_{\mathrm{alloc}}^\star \ge 0.5$, and stage 2 removes them: the best
available allocation is declared unconditionally, whatever its probability.
\sestina{} removes the mechanism that used to raise pressure in ordinary play, and
\S\ref{sec:agent-stall} is what now supplies stage 2.

\subsection{Urgency-escalation removal}
\label{sec:agent-urgency}

\vsix{} inherited a declaration-urgency escalation keyed to an event counter. Past a threshold of 220
public events the declaration rule lowers its team-ownership floor and permits declarations with
estimated correctness probabilities near 0.5, which costs \vsevenCliff{}~pp to a policy held in that
regime. No adversary
reaches it, since the longest game any adversarial search produced against an unmodified opponent is
\vsevenCliffIncumbentTail{} events, but the machinery is defective in ordinary play as well:
\S\ref{sec:development-alloc} reports that removing it eliminates \vsevenLoneUrgencyShare\% of the
entire declaration allocation-error class. \sestina{} disables it
(\texttt{pool}, \texttt{oppfloor}, \texttt{force}, \texttt{askfloor}).

The escalation was, however, also the policy's termination safeguard, so removing this mechanism
requires an alternative stall safeguard.

\subsection{Deduction-state stall detection}
\label{sec:agent-stall}

The escalation \sestina{} removes was keyed to a clock. The condition the clock stood in for is
that the monitored deduction state has not changed, and that condition is directly observable from
the seat's own deduction state.

Let $\chi_t = \mathrm{hash}\bigl(\mathrm{own}_t(\cdot),\, \mathcal{M}_t(\cdot),\, U_t,\,
(\text{hand sizes})_t,\, (\text{live half-suits})_t\bigr)$, the digest of the hard certificates the
seat holds. It deliberately \emph{excludes} the list of outstanding (C5) disjunctions: a repeated ask
appends a duplicate disjunction carrying no information, so counting that list as progress would
blind the detector to the one loop it exists to break. Define the
no-progress counter
\begin{equation}
  \nu_t \;=\;
  \begin{cases}
    0, & \chi_t \neq \chi_{t-1},\\
    \nu_{t-1} + 1, & \chi_t = \chi_{t-1},
  \end{cases}
  \label{eq:stall}
\end{equation}
and escalate to pressure stage 1 when $\nu_t \ge K$ and to stage 2 when $\nu_t \ge 2K$, with
$K = \texttt{\vsevenOptStallK{}}$ and $2K = \texttt{\vsevenOptStallTwo{}}$ in \sestina{}.

\paragraph{Scope of the termination safeguard.} Consider
$\Phi_t = \sum_{c \in U_t}\bigl(|\mathcal{M}_t(c)| - 1\bigr) + |U_t| + |\mathcal{H}_t|$. Certificates
only accumulate: a seat excluded from a card's candidate set is never restored, a resolved card is
never unresolved again, and a claimed half-suit never returns. Every event that changes $\chi$ by one
of those routes therefore strictly decreases $\Phi_t$, which is a non-negative integer, so there can be
only finitely many of them. This potential-function argument does not establish termination,
because transfers can change the monitored state without decreasing the proposed potential: the
digest also contains hand sizes, so a successful ask always counts as progress, and a card moving
between two seats is not a monotone gain; nothing above forbids a card being won back later. The rule therefore bounds the mode
of non-termination that actually occurs, a run of asks producing no new certificate, and not every
mode. That is why this paper calls it a safeguard, and why \S\ref{sec:discussion-engineering} records
termination as empirical.

The threshold was set from measurement. Over \vsevenStallMirrorGames{} mirror games \vsix{} never
produces a no-progress run longer than \vsevenStallIncumbentRun{}, while the deliberately deadlocked
negative control of \S\ref{sec:evaluation-gate} produces \vsevenStallFrozenRun{}. Any
$K \in [12,60]$ was therefore not reached in the sampled ordinary games and was reached in the
tested deadlock control, after $K$ and then $2K$ no-progress events. Armed, the rule's mirror
transcript digest is byte identical to the same configuration without it: in those mirror games its
effect on play is identically zero, because $\nu_t$ never reaches $K$.
\S\ref{sec:results-attribution} confirms that on holdout material. Note also that $\nu_t$ is a
function of the seat's own hand and the public record alone, so the rule opens no information
channel.

\subsection{Guarded determinized search}
\label{sec:agent-search}

\sestina{} enables the search \vsix{} built, measured, and then deployed switched off. It runs only
in the endgame, when at most \texttt{\vsevenOptMaxq{}} cards remain unresolved, and it revises the
blueprint's choice only when a standard-error-penalised paired comparison favours the revision.

\paragraph{Deal sampling.} The search draws $D_{\mathrm{det}} = \texttt{\vsevenOptDet{}}$
determinizations from the constraint-conditioned distribution \eqref{eq:posterior}, using a
per-card capacity dynamic program as a uniform sampler over deals satisfying (C1)--(C4), with the
(C5) disjunctions enforced by rejection, so an accepted draw is an exact sample from
\eqref{eq:posterior}. This sampler is a distinct object from the block program of \S\ref{sec:inference-factor},
which enumerates (C5) inside its blocks instead. $D_{\mathrm{det}}$ is a target rather than a
guarantee: draws are attempted up to a cap and the rule proceeds with however many survive
rejection. If no draw survives, or the sampler cannot be built, the decision falls back to the
\vfive{} scoring path rather than to the blueprint's own choice; that path scores with the
twenty-feature vector alone and breaks ties in enumeration order. No artifact of this cycle
records how often the fallback fired for the frozen configuration. The behavioural weighting of \eqref{eq:prior} is then applied as an
importance weight rather than baked into the sampler, which keeps the sampler exact and the
weighting auditable:
\begin{equation}
  \tilde w_d \;=\; \prod_{c \in U_t} \exp\Bigl(\operatorname{clip}\bigl(\theta\, a_{o_d(c),H(c)}
  - \varphi\,(A_{o_d(c)} - a_{o_d(c),H(c)}),\, \pm 2.6\bigr)\Bigr),
  \qquad
  w_d = \frac{\tilde w_d}{\sum_{d'} \tilde w_{d'}},
  \label{eq:iw}
\end{equation}
where $o_d(c)$ is the seat holding card $c$ in world $d$.

\paragraph{Rollout evaluation.} The candidate set is the top $\texttt{\vsevenOptCand{}}$ blueprint
candidates by score, extended through any tie at the cut up to a cap of twenty, with index $0$
reserved for the blueprint's own choice $a_{(1)}$, the first candidate after the score sort rather
than the output of \eqref{eq:tiebreak}. For each candidate $r$ and world $d$, all six seats are reconstructed \emph{at their own information
sets} under $d$, each playing a cheaper blueprint than the acting agent's own: the \vfive{} policy
with an independent-marginal belief and no re-scoring pass and the game is played to a terminal state or until \texttt{\vsevenOptDepth{}} further public events
have elapsed. A run reaching the end returns the acting team's half-suit differential; a run cut at the limit
returns a \emph{material} evaluation of the position instead, being the scored half-suit differential
plus the team's expected control of the live half-suits. That is \vsix{}'s own hand-written leaf, not
a learned one: \S\ref{sec:development-leaf} reports what happened when a fitted replacement was
built. Write $V_{r,d}$ for either. This is the point at which the method departs from standard
determinization: the continuation players are not made clairvoyant. \S\ref{sec:game-why} explains why
that matters here, namely that the recorded omniscient-play probe (\S\ref{sec:game-why})
motivates concern that a clairvoyant continuation model would discriminate weakly between candidate
actions; the probe does not establish action-level values. All candidates share the same $D_{\mathrm{det}}$ worlds, so the comparison is paired and the
noise that matters is the noise in the difference.

\paragraph{Guarded action selection.} For each candidate $r \ge 1$ form the paired difference
against the blueprint,
\begin{equation}
  m_r \;=\; \sum_{d} w_d\,\bigl(V_{r,d} - V_{0,d}\bigr),
  \qquad
  \mathrm{var}_r \;=\; \sum_{d} w_d\,\bigl(V_{r,d} - V_{0,d} - m_r\bigr)^2,
  \label{eq:paired}
\end{equation}
and with the Kish effective sample size $n_{\mathrm{eff}} = 1/\sum_d w_d^2$ take
\begin{equation}
  \mathrm{se}_r = \sqrt{\frac{\mathrm{var}_r}{n_{\mathrm{eff}} - 1}},
  \qquad
  \mathrm{LCB}_r = m_r - \kappa_r\,\mathrm{se}_r + \eta\,\bigl(u(a_{(r+1)}) - u(a_{(1)})\bigr).
  \label{eq:lcb}
\end{equation}
The search deviates to $\arg\max_r \mathrm{LCB}_r$ if that maximum is strictly positive, and
otherwise plays the blueprint's choice. When the effective sample size is at most $1.5$, which
includes the case of a single surviving world or one dominant weight, $\mathrm{se}_r$ is set to an
effectively infinite value, so no candidate outside the tied set can deviate and within it the raw
weighted paired mean decides. The shrinkage is $\kappa = \texttt{\vsevenOptKappa{}}$. The term $\eta$ blends the blueprint's own
ranking back in and is \textbf{zero in this configuration}, so \eqref{eq:lcb} reduces to a pure
paired lower confidence bound.

Equation~\eqref{eq:lcb} is the action-selection guard used in the evaluated configuration. Chosen
instead by an unguarded $\arg\max_r m_r$, the same search is \vsixCurseGap{}~pp \emph{worse} than
the blueprint it searches from (\S\ref{sec:background-vsix}): one world's return is a final
half-suit differential whose spread is an order of magnitude larger than genuine differences between
candidates, so the maximum selects on residual noise, and does so more confidently the more
candidates it is offered. Requiring $\kappa$ standard errors of a \emph{paired} difference removed
that measured deficit in the \vsix{} experiment. The rule is a lower-confidence-bound heuristic
rather than a test with established finite-sample coverage: at most $D_{\mathrm{det}}$ weighted
worlds enter it, several candidates are compared against the same reference, and the effective
sample size can be small. What the record establishes is that the guard reduced the measured
deficit, not that it guarantees an improvement over the blueprint.

One detail is worth stating because \S\ref{sec:development-tie} turns on it. Within the tied set
$\mathcal{E}_t$ of \eqref{eq:tieset} the rule applies $\kappa_r = 0$, that is, no shrinkage at all,
which is precisely where its selection signal is most noise-dominated.

\subsection{Coordination and information constraints}
\label{sec:agent-coordination}

\sestina{} deploys as three identical agents. They hold different private hands, they may not
signal, and the only channels the rule variant leaves are the forced-endgame willingness bit and the
successor a cardless turn-holder names (\S\ref{sec:game-dialect}), for both of which
\S\ref{sec:game-dialect} reports no resolved effect from the tested interventions. The
coordination assumptions are therefore the following.

They rest on determinism plus a shared observation. Consider first the idealised non-searching
policy, and let $\Pi$ denote it: a map from $(\text{own hand}, \text{public record})$ to an action.
Every teammate runs the same $\Pi$ and observes the same public record $\mathcal{P}_t$. Teammate $j$ cannot compute $\Pi$'s output at seat
$i$, because it does not know $i$'s hand. But it knows $\Pi$ itself, and it knows $\mathcal{P}_t$, so
it can compute the \emph{distribution} of $i$'s action induced by its own posterior over $i$'s hand,
\begin{equation}
  \Pr\bigl[\,a_i = a \mid \mathcal{I}_j^t\,\bigr]
  \;=\; \sum_{D \,\models\, \mathcal{I}_j^t}
  \Pr\bigl[D \mid \mathcal{I}_j^t\bigr]\;
  \ind\!\left[\Pi\bigl(\text{hand}_i(D),\, \mathcal{P}_t\bigr) = a\right],
  \label{eq:coord}
\end{equation}
and, having observed $a_i$, condition on it. That is exactly the (C5) reasoning of
\S\ref{sec:inference-constraints} carried one step further: a legality certificate says the asker
held \emph{some} card of the half-suit, while \eqref{eq:coord} says the asker held a hand on which
$\Pi$ preferred that ask. The deployed approximation does not evaluate \eqref{eq:coord} in full,
which would require enumerating hands; it captures the first-order part of it through the ask-tally
prior \eqref{eq:prior}. The white-box adversary of \S\ref{sec:results-whitebox} is the construction
that approximates \eqref{eq:coord} against \sestina{} rather than for it: it evaluates the one-step
action likelihood on exact samples and accumulates per-card evidence additively rather than
carrying the joint conditioning across the game, and \S\ref{sec:development-instrument} reports
what that inversion is worth.

Two conditions make this work, and both are engineering constraints rather than assumptions. $\Pi$
must be \emph{specified}: a teammate reasoning about seat $i$ must be able to evaluate
$\Pi(\cdot,\mathcal{P}_t)$ on a hypothesised hand from the rule alone. An unstable sort does not
break determinism, since it is deterministic for a fixed input and build, but its result is fixed by
an implementation detail rather than by the policy specification, so a teammate would have to model
the implementation; \eqref{eq:tiebreak} replaces it with a specified rule on the non-search path.
And $\Pi$ must be common knowledge, which for three copies of one published agent it is. No shared
secret and no pre-play randomisation are used, which is the sense in which the team's coordination
is public (\S\ref{sec:threat}).

Equation~\eqref{eq:coord} is stated for the idealised non-searching policy. The deployed agent
departs from it in two ways. When the search of \S\ref{sec:agent-search} fires, the action also
depends on the seat's reset seed, which is private to the seat, so the deployed policy is a function
of $(\text{own hand}, \mathcal{P}_t, \text{reset seed})$ on that path (\S\ref{sec:agent-tiebreak});
a teammate can still compute the distribution of the action given a hypothesised hand only by
modelling that randomness. And the cross-deal state persistence described next introduces a further
dependence that neither a teammate nor \eqref{eq:coord} models.

\paragraph{Cross-deal state persistence.} Agents are constructed once per worker thread and
reused across the deals that thread is given, and some per-agent state survives the inter-deal reset,
reachable only under the search. This violates the premise that $\Pi$ depends on
$(\text{own hand}, \mathcal{P}_t)$ alone: a decision can depend on which deals the thread happened to
process earlier. The carrier was identified as the rollout engine, which the reset never touches; the
specific field was not. It was deliberately not repaired, because the one-line fix changes the play of
every searching configuration by roughly one deal in \vsevenSsixIncidence{}, and applying it after the
configuration was frozen would defeat the purpose of freezing.
\S\ref{sec:results-controls} measures its magnitude on holdout material, and
\S\ref{sec:discussion-engineering} states what it costs the argument above.

\subsection{Decision procedure}
\label{sec:agent-full}

Collecting the pieces, one turn of \sestina{} at seat $s$ and time $t$ runs as follows. The steps
describe the deployed approximate path, not the reference objects of \S\ref{sec:inference}.

\begin{enumerate}
\item \textbf{Update beliefs.} Fold the new public events into the candidate sets
  $\mathcal{M}_t(\cdot)$ by (C1)--(C5), then compute $\hat\mu$ by \eqref{eq:prior} and
  \eqref{eq:sinkhorn}. Update the no-progress counter $\nu_t$ by \eqref{eq:stall}.
\item \textbf{Consider declaring.} For each live half-suit passing the cheap ownership gates,
  enumerate the capacity-feasible allocations of its unresolved cards to the acting team and select
  the allocation maximising the product of approximate card marginals. Estimate its probability
  $\hat\pi_{\mathrm{alloc}}^\star$ by sequential conditioning and belief refitting. If the
  team-ownership estimate clears its floor, declare when $\Delta > \delta$ in \eqref{eq:declev};
  declare regardless if a pressure override fires, and answer the forced-endgame ladder by its
  thresholds (\S\ref{sec:agent-declare}).
\item \textbf{Otherwise ask.} Enumerate the legal set $\mathcal{A}_t$ by \eqref{eq:legal} and apply
  the live-ask filter of \S\ref{sec:agent-policy}, removing candidates whose target is excluded by a
  certificate; if the filter removes every candidate, restore the legal set. Score the resulting
  candidates by \eqref{eq:askscore}, including the contestation term \eqref{eq:contest}.
\item \textbf{Resolve ties.} Form $\mathcal{E}_t$ by \eqref{eq:tieset}. If the endgame condition does
  not hold, return \eqref{eq:tiebreak} and stop.
\item \textbf{Search.} Otherwise sample $D_{\mathrm{det}}$ worlds with weights \eqref{eq:iw}, roll out
  each candidate in each world, and return the candidate selected by \eqref{eq:lcb}. This is the
  blueprint's own choice, the first candidate after the score sort, unless a candidate's paired
  lower confidence bound is positive; the bound uses $\kappa_r = \kappa$ standard errors for
  candidates outside the tied set $\mathcal{E}_t$ and $\kappa_r = 0$ within it, and if no world
  survives rejection the \vfive{} scoring path decides (\S\ref{sec:agent-search}).
\end{enumerate}

Steps 1--4 read only the acting seat's own hand and the public record; step 5 also reads the seat's
own reset seed. \S\ref{sec:agent-coordination} records that, in the evaluated implementation, a
searching agent's decisions can additionally depend on state carried across deals, so the property
that every decision is a function of the listed information holds for the non-searching path and is
an approximation for the searching one. The side-channel audit of \S\ref{sec:evaluation-side} tests
the property mechanically rather than by inspection, and \S\ref{sec:results-controls} reports where
it fails.

